% The Tensor Decomposition of the Shielding Perturbation: 723 Mutants
% Standalone section for \input{} into thesis chapter.
% Requires: booktabs, siunitx, amsmath

\section{The Tensor Decomposition of the Shielding Perturbation: 723 Mutants}
\label{sec:tensor-decomposition}

The NMR shielding extractor computes 121 geometric kernels---rank-2
spherical tensors ($L=2$, five components)---at each atom in a protein,
decomposing the shielding perturbation from aromatic ring deletion into
its physical sources.  Before using these kernels for calibration, we
verify that they reproduce 28 known properties of the shielding
interaction, drawn from the NMR and physics literature.

All measurements are on WT-ALA mechanical mutants: backbone coordinates
fixed, aromatic sidechain replaced with alanine.  The DFT shielding
tensor difference $\Delta\sigma_{T_2}$ is the reference throughout.
The dataset comprises 110 proteins and 69{,}080 matched atoms.


%%% ---------------------------------------------------------------
\subsection{Distance Dependence}
\label{ssec:distance}

The multipole expansion determines the radial falloff of each kernel
family.  A magnetic dipole (ring current) falls off as $r^{-3}$
\cite{johnson1958}; an electric quadrupole as $r^{-5}$
\cite{stone2013}.  Log-log regression on 58{,}193 atom--ring pairs
($r > \SI{2}{\angstrom}$) gives the slopes listed in
Table~\ref{tab:distance-slopes}.

\begin{table}[htbp]
  \centering
  \caption{Distance falloff of kernel $T_2$ magnitude.  Expected slopes
    from multipolar theory; measured slopes from log-log regression.}
  \label{tab:distance-slopes}
  \begin{tabular}{llcc}
    \toprule
    Kernel & Physical source & Expected & Measured \\
    \midrule
    Biot--Savart    & Magnetic dipole (ring current) & $-3$ & $-3.04$ \\
    Haigh--Mallion  & Surface integral (ring current) & $-3$ & $-3.06$ \\
    Pi-Quadrupole   & Electric quadrupole             & $-5$ & $-5.05$ \\
    \bottomrule
  \end{tabular}
\end{table}

Three-significant-figure agreement with multipolar theory.  The
distance dependence is not fitted---it is encoded in the geometry of
the kernel computation (Biot--Savart law \cite{johnson1958},
Stone $T$-tensor formalism \cite{stone2013}).
See Figure~\ref{fig:distance_falloff}.


%%% ---------------------------------------------------------------
\subsection{Angular Dependence}
\label{ssec:angular}

The Biot--Savart kernel magnitude depends on the polar angle $\theta$
relative to the ring normal.  The $(1 - 3\cos^2\theta)$ angular pattern
of a magnetic dipole \cite{pople1956} predicts axial enhancement, which
is observed directly (Figure~\ref{fig:angular_peaking}).

\begin{table}[htbp]
  \centering
  \caption{Biot--Savart $T_2$ magnitude by polar angle bin.}
  \label{tab:angular-bins}
  \begin{tabular}{ccc}
    \toprule
    Angle range & Mean BS magnitude & Relative \\
    \midrule
    \SIrange{0}{17}{\degree}  & 0.0385 & 1.00 \\
    \SIrange{17}{34}{\degree} & 0.0331 & 0.86 \\
    \SIrange{34}{52}{\degree} & 0.0237 & 0.62 \\
    \SIrange{52}{69}{\degree} & 0.0172 & 0.45 \\
    \SIrange{69}{90}{\degree} & 0.0139 & 0.36 \\
    \bottomrule
  \end{tabular}
\end{table}

Monotonic decrease with a $2.8\times$ axial-to-equatorial ratio,
consistent with $\lvert 1 - 3\cos^2\theta \rvert$ over the
corresponding angular ranges \cite{case1995}.


%%% ---------------------------------------------------------------
\subsection{Element Dependence}
\label{ssec:element}

The literature predicts that ring current dominates proton shielding
\cite{boyd2002} while the electric field gradient dominates heavy-atom
shielding \cite{sahakyan2013}.  These predictions concern isotropic
shifts; we test them in the $T_2$ (anisotropic) channel.

\begin{table}[htbp]
  \centering
  \caption{Ridge $R^2$ by kernel family and element.  Bold entries mark
    the dominant mechanism for each element.}
  \label{tab:element-r2}
  \begin{tabular}{lrcccc}
    \toprule
    Element & $n$ & Ring current & EFG & Bond aniso & All \\
    \midrule
    H & 35{,}080 & \textbf{0.909} & 0.770 & 0.139 & 0.935 \\
    C & 20{,}673 & 0.156 & \textbf{0.431} & 0.041 & 0.546 \\
    N &  6{,}100 & \textbf{0.325} & 0.100 & 0.070 & 0.434 \\
    O &  6{,}927 & \textbf{0.317} & 0.196 & 0.055 & 0.387 \\
    \midrule
    Pooled & 69{,}080 & 0.213 & 0.318 & 0.027 & 0.385 \\
    \bottomrule
  \end{tabular}
\end{table}

The pooled row suggests EFG dominance---a composition artefact.
Per-element analysis reveals that H atoms respond predominantly to ring
current ($R^2 = 0.909$) while C atoms respond to the electric field
gradient ($R^2 = 0.431$).  N and O show multi-mechanism behaviour
(Section~\ref{ssec:nitrogen}).  The weighted per-element $R^2$ with all
kernels is 0.718, exceeding the pooled value by 0.333
(Figure~\ref{fig:element_r2}).

The tensor alignment (cosine similarity between the EFG kernel $T_2$
and the DFT target $T_2$) confirms the element ordering without any
fitting: mean $\lvert\cos\rvert = 0.930$ for H, $0.863$ for C, $0.790$
for N, $0.737$ for O (Figure~\ref{fig:efg_cos_by_element}).  Forward
selection results are shown in Figure~\ref{fig:forward_selection}.


%%% ---------------------------------------------------------------
\subsection{Nitrogen}
\label{ssec:nitrogen}

Nitrogen does not fit the ring-current-versus-EFG dichotomy.  No single
kernel family dominates: the largest single-kernel $R^2$ is 0.080, and
five mechanisms each contribute 0.015--0.080 in forward selection (EFG,
pi-quadrupole, Biot--Savart, McConnell bond anisotropy, ring
susceptibility).

This is consistent with four observations from the literature:

\begin{enumerate}
  \item ${}^{15}$N CSA is large (${\sim}\SI{166}{ppm}$) and
    site-variable ($\pm\SI{20}{ppm}$), dominated by backbone angles and
    hydrogen bonding \cite{yao2010,hall2006}.  In mechanical mutants
    these are held constant; the residual comes from secondary
    mechanisms.

  \item ${}^{15}$N shielding is paramagnetic-term-dominated
    \cite{saito2010}, making it sensitive to multiple perturbation
    channels simultaneously.

  \item The isotropic ring current effect on ${}^{15}$N is negligible
    ($<0.6\%$ of variance) \cite{han2011}, but this does not preclude
    an anisotropic ($T_2$) effect.  Hall \& Fushman \cite{hall2006}
    noted aromatic residues as ${}^{15}$N CSA outliers in protein~GB3.

  \item MopacMC entering the N forward selection at rank~4 reflects
    electronic coupling through the peptide bond (partial double-bond
    character, Wiberg bond order ${\sim}1.3$--$1.5$).
\end{enumerate}


%%% ---------------------------------------------------------------
\subsection{Model Agreement: Biot--Savart and Haigh--Mallion}
\label{ssec:bs-hm}

The Biot--Savart \cite{johnson1958} and Haigh--Mallion
\cite{haighmallion1979} models compute the same physical effect---the
ring current magnetic field---with different mathematics (wire-loop
segments versus surface integral).  Case \cite{case1995} found no
significant difference between them for isotropic shifts.  We extend
the comparison to the $T_2$ channel.

\begin{table}[htbp]
  \centering
  \caption{Biot--Savart versus Haigh--Mallion agreement in the $T_2$
    channel.}
  \label{tab:bs-hm}
  \begin{tabular}{lcc}
    \toprule
    Property & Value & Agreement \\
    \midrule
    Distance slope (BS / HM) & $-3.04$ / $-3.06$ & identical \\
    $T_2$ magnitude ratio (per atom) & $1.009 \pm 0.037$ & $1\%$ \\
    \bottomrule
  \end{tabular}
\end{table}

The $T_2$ angular agreement (cosine similarity) is distance-dependent:
$\langle\lvert\cos\rvert\rangle = 0.9977$ in the near field
($r < \SI{5}{\angstrom}$) and $0.9999$ in the far field
($r > \SI{8}{\angstrom}$).  Both models converge to the magnetic dipole
limit at long range; the 0.2\% near-field disagreement reflects the
difference between wire-loop and surface geometry at close range
(Figure~\ref{fig:bs_hm_ratio}, Figure~\ref{fig:bs_hm_convergence}).


%%% ---------------------------------------------------------------
\subsection{Ring-Type Dependence}
\label{ssec:ring-type}

Each aromatic ring type is fitted independently with its four
ring-specific kernels (Biot--Savart, Haigh--Mallion, EFG, susceptibility).

\begin{table}[htbp]
  \centering
  \caption{Self-fit $R^2$ by ring type (ring-specific kernels only).}
  \label{tab:ring-type-r2}
  \begin{tabular}{lc}
    \toprule
    Ring type & $R^2$ \\
    \midrule
    TRP (perimeter)  & 0.125 \\
    TRP (benzene)    & 0.114 \\
    TYR              & 0.106 \\
    PHE              & 0.098 \\
    TRP (pyrrole)    & 0.095 \\
    HIE (imidazole)  & 0.062 \\
    \bottomrule
  \end{tabular}
\end{table}

The PHE/TYR $R^2$ ratio is 0.84, consistent with their similar ring
current intensities ($I_{\text{PHE}} = 1.46$, $I_{\text{TYR}} = 1.24$
in the Haigh--Mallion model \cite{case1995}).

HIE is worst: the circular-loop Biot--Savart model fails for the
asymmetric five-membered ring with two nitrogen positions.  Ring current
kernels alone are never sufficient ($R^2 = 0.06$--$0.13$); they require
the EFG and other kernel families.


%%% ---------------------------------------------------------------
\subsection{Mechanical Mutant Controls}
\label{ssec:controls}

Two negative controls verify that the kernels respond to the aromatic
ring deletion, not to the unchanged backbone.

\begin{table}[htbp]
  \centering
  \caption{Negative control kernels.  Both are consistent with zero
    response to the mechanical mutation.}
  \label{tab:controls}
  \begin{tabular}{llcc}
    \toprule
    Test & Kernel & $R^2$ & Expected \\
    \midrule
    H-bond unchanged   & HBond total    & 0.002 & ${\approx}0$ \\
    Solvation minor    & DeltaAPBS EFG  & 0.005 & $\ll$ EFG$_{\text{aro}}$ \\
    \bottomrule
  \end{tabular}
\end{table}

Hydrogen-bond shielding (backbone N--H$\cdots$O=C) does not change in a
mechanical mutant; the kernel correctly returns zero.  The solvation
correction (APBS continuum minus vacuum Coulomb) is 2\% of the direct
aromatic EFG ($R^2 = 0.230$), confirming that the aromatic shielding
perturbation is a local, through-space effect.


%%% ---------------------------------------------------------------
\subsection{Inter-Kernel Geometry}
\label{ssec:inter-kernel}

\paragraph{Three effective dimensions.}
The 91-kernel $T_2$ space has three effective dimensions at 98\%
variance (eigenspectrum analysis), corresponding to three physical
source geometries: (i)~the ring current magnetic dipole (Biot--Savart,
Haigh--Mallion, ring susceptibility); (ii)~the electric field gradient
(Coulomb EFG, MOPAC EFG); and (iii)~bond magnetic anisotropy
(McConnell).

\paragraph{BS--EFG angular independence.}
The mean $\lvert\cos(\text{BS}, \text{EFG})\rvert = 0.684$ across
58{,}193 atom--ring pairs.  This lies between collinear (1.0) and
random in five dimensions (0.45): the two kernel families share
geometric information (both depend on the ring position) but encode
different angular patterns (magnetic dipole versus electric quadrupole).

\paragraph{BS additivity for fused rings.}
For tryptophan's fused ring system,
$T_2(\text{TRP}_5) + T_2(\text{TRP}_6) = T_2(\text{TRP}_9)$ at machine
precision.  The wire-segment Biot--Savart law is exactly additive for
non-overlapping current loops.


%%% ---------------------------------------------------------------
\subsection{Limitations}
\label{ssec:limitations}

\begin{enumerate}
  \item \textbf{Per-protein ceiling.}  Per-protein ridge $R^2 = 0.81$;
    global (cross-protein) ridge $R^2 = 0.35$; MLP $R^2 = 0.61$.  The
    gap represents per-protein structural variation (bulk electrostatic
    environment, protein shape) that local geometric kernels do not
    capture.

  \item \textbf{Near-field accuracy.}  The Biot--Savart and
    Haigh--Mallion models approximate the true $\pi$-electron current
    distribution.  At $r < \SI{4}{\angstrom}$ the wire-loop and
    surface-integral models diverge from the quantum-mechanical
    current density.  Per-ring kernels partially address this, but the
    functional form is fixed.

  \item \textbf{HIE modelling.}  The worst ring type (self-fit
    $R^2 = 0.062$).  The symmetric circular-loop approximation fails
    for the asymmetric five-membered ring.

  \item \textbf{Isotropic gap.}  These kernels describe the $T_2$
    (anisotropic) shielding perturbation.  The $T_0$ (isotropic shift)
    requires the Buckingham $A$ coefficient \cite{buckingham1960} and
    ring current intensity as additional parameters, fitted separately.

  \item \textbf{No structural relaxation.}  Mechanical mutants hold the
    backbone fixed.  Real mutations cause structural rearrangement that
    changes the kernel geometry.  The calibrated kernels apply to a
    fixed structure, not to a mutation experiment.
\end{enumerate}
